% Timing diagram of the interrupt request port of the pipelined core
% (MICROARCH.md §6), rendered to interrupt.png by "make diagrams" in the
% top-level Makefile. Two IRQs through vector 9:
%
% (a) A request that comes while the next instruction has not begun. It is
%     taken at once, before that instruction, which goes with the flush.
% (b) A request that comes just after the next instruction has issued its
%     data read. That instruction is committed, so the request waits until
%     it has retired, and is taken before the one after.
%
% The port is valid/ready with a vector, as in qnice_cpu: the device holds
% irq_valid_i and irq_vector_i until irq_ready_o, and irq_ready_o is
% combinational in irq_valid_i and WRITE's registers. Taking the request
% flushes the pipeline (flush_o), and WRITE reads the vector table entry
% at VBASE * 256 + 9 = 0x0009 on the data bus (VEC). The read goes into
% DBUS's register, which drives the bus from the next cycle on
% (MICROARCH.md §5.3), so it is on the bus a cycle after the take. The entry is reported
% (retire_o, "IRQ") in the cycle after its data comes, as WRITE takes the
% vector from DBUS's FIFO (MICROARCH.md §5.3), and FETCH is redirected
% there (target_o) in the next cycle, from a register (MICROARCH.md §5.4).
% FETCH is redirected in the cycle after the take too, and then stops
% until the entry is done. The target of the take is the
% address of the instruction after the dropped one (u_v.inst.next_pc in
% pipe_write.vhd); FETCH never fetches from it, as it is redirected again
% to the entry.
%
% All values below are read off a GHDL simulation of hdl/test/irqwave.nasm
% on the pipelined core, with "--irq 2:9 --irq 10:9" (make test,
% INTR_irqwave) and a memory that never stalls. Cycle 0 is the first cycle
% after reset. Panel (a) starts at cycle 9 and panel (b) at cycle 36.
% Changing the program or the core invalidates them. The addresses:
%
%   0010-0024  the main program: 0010 sets E, 0012-001A and 001E-0022 are
%              ADD R1, R1, 1, 001C is MOVE R2, @data, and 0024 is HALT
%   0025, 0027 the routine: ADD R3, R3, 1, then RTI
%   0028       data, which PREPARE briefly holds as an instruction after RTI
%
% A bus is drawn gray (x) in every cycle in which it carries nothing:
% irq_vector_i while irq_valid_i is low, prep_r and w_r.uop while they hold
% no micro-operation, dbus_addr_o while dbus_stb_o is low, target_o
% while flush_o is low, and retire_o while it reports no step. prep_r and
% w_r.uop show the address of their micro-operation's instruction.
%
% (a) t=2   The request comes. WRITE holds the first micro-operation of
%           0012, which has not begun, so it takes the request at once.
%     t=3   The vector read goes into DBUS's register; t=4 it is on the
%           bus; t=5 its data; t=6 the entry.
%     t=16  The routine's RTI retires and flushes back to 0012.
% (b) t=1   PREPARE loads the read of 001C, while WRITE retires 001A.
%     t=2   The request comes. 001C has issued its read (commits = 1), so
%           no request may be taken before it.
%     t=3   001C retires with its data; t=4 the request is taken before
%           001E; t=8 the entry.

\documentclass[varwidth = 40cm, border = 8pt]{standalone}
\usepackage{timing}

% Four-digit addresses fit a one-cycle box at this scale.
\renewcommand{\wavescale}{1.2}

\begin{document}

{\sffamily (a) A request taken at once}\par\medskip

\begin{wave}{15}{18}
 \nextwave{irq\_valid\_i}            \bit{0}{2} \bit{1}{1} \bit{0}{15}
 \nextwave{irq\_vector\_i}           \unknown[x]{2} \known{09}{1} \unknown[x]{15}
 \nextwave{irq\_ready\_o}            \bit{0}{2} \bit{1}{1} \bit{0}{15}
 \nextwave{PREPARE: prep\_r}         \known{0010}{1} \known{0012}{1} \known{0014}{1} \unknown[x]{11} \known{0025}{1} \known{0027}{1} \known{0028}{1} \unknown[x]{1}
 \nextwave{PREPARE: load\_o}         \bit{0}{18}
 \nextwave{WRITE: w\_r.uop}          \unknown[x]{1} \known{0010}{1} \known{0012}{1} \unknown[x]{12} \known{0025}{1} \known{0027}{1} \unknown[x]{1}
 \nextwave{WRITE: w\_r.mode}         \known{RUN}{3} \known{VEC}{4} \known{RUN}{11}
 \nextwave{WRITE: w\_r.commits}      \known{0}{18}
 \nextwave{WRITE: w\_r.in\_service}  \bit{0}{7} \bit{1}{10} \bit{0}{1}
 \nextwave{dbus\_stb\_o}             \bit{0}{4} \bit{1}{1} \bit{0}{13}
 \nextwave{dbus\_addr\_o}            \unknown[x]{4} \known{0009}{1} \unknown[x]{13}
 \nextwave{dbus\_ack\_i}             \bit{0}{5} \bit{1}{1} \bit{0}{12}
 \nextwave{WRITE: flush\_o}          \bit{0}{2} \bit{1}{1} \bit{0}{3} \bit{1}{1} \bit{0}{9} \bit{1}{1} \bit{0}{1}
 \nextwave{WRITE: target\_o}         \unknown[x]{2} \known{0014}{1} \unknown[x]{3} \known{0025}{1} \unknown[x]{9} \known{0012}{1} \unknown[x]{1}
 \nextwave{retire\_o}                \unknown[x]{1} \known{0010}{1} \unknown[x]{4} \known{IRQ}{1} \unknown[x]{8} \known{0025}{1} \known{0027}{1} \unknown[x]{1}
\end{wave}

\bigskip
{\sffamily (b) A request that waits for an instruction that has issued its data read}\par\medskip

\begin{wave}{15}{10}
 \nextwave{irq\_valid\_i}            \bit{0}{2} \bit{1}{3} \bit{0}{5}
 \nextwave{irq\_vector\_i}           \unknown[x]{2} \known{09}{3} \unknown[x]{5}
 \nextwave{irq\_ready\_o}            \bit{0}{4} \bit{1}{1} \bit{0}{5}
 \nextwave{PREPARE: prep\_r}         \known{001A}{1} \known{001C}{1} \known{001E}{2} \known{0020}{1} \unknown[x]{5}
 \nextwave{PREPARE: load\_o}         \bit{0}{1} \bit{1}{1} \bit{0}{8}
 \nextwave{WRITE: w\_r.uop}          \known{0018}{1} \known{001A}{1} \known{001C}{2} \known{001E}{1} \unknown[x]{5}
 \nextwave{WRITE: w\_r.mode}         \known{RUN}{5} \known{VEC}{4} \known{RUN}{1}
 \nextwave{WRITE: w\_r.commits}      \known{0}{2} \known{1}{2} \known{0}{6}
 \nextwave{WRITE: w\_r.in\_service}  \bit{0}{9} \bit{1}{1}
 \nextwave{dbus\_stb\_o}             \bit{0}{2} \bit{1}{1} \bit{0}{3} \bit{1}{1} \bit{0}{3}
 \nextwave{dbus\_addr\_o}            \unknown[x]{2} \known{0028}{1} \unknown[x]{3} \known{0009}{1} \unknown[x]{3}
 \nextwave{dbus\_ack\_i}             \bit{0}{3} \bit{1}{1} \bit{0}{3} \bit{1}{1} \bit{0}{2}
 \nextwave{WRITE: flush\_o}          \bit{0}{4} \bit{1}{1} \bit{0}{3} \bit{1}{1} \bit{0}{1}
 \nextwave{WRITE: target\_o}         \unknown[x]{4} \known{0020}{1} \unknown[x]{3} \known{0025}{1} \unknown[x]{1}
 \nextwave{retire\_o}                \known{0018}{1} \known{001A}{1} \unknown[x]{1} \known{001C}{1} \unknown[x]{4} \known{IRQ}{1} \unknown[x]{1}
\end{wave}

\end{document}
